\chapter{Problem Identification and Formulation}
\label{ch:problem}

\section{Background and Motivation}

Remote onboarding and account recovery frequently compare a selfie or short face
video with identity evidence. The comparison is useful only if the media came
from the person and capture session being verified. An attacker may instead
present a printed photograph, replay a video on another screen, wear a mask,
display a manipulated video, or bypass the physical camera and inject forged
media. ENISA consequently separates camera presentation from injection attacks
and recommends layered controls spanning the client and server
\cite{enisa2022,enisa2024}. NIST likewise treats sensor authenticity, protected
channels, presentation-attack detection (PAD), and forged-media analysis as
complementary controls rather than substitutes \cite{nist80063a4}.

This problem is especially consequential in financial verification. A false
acceptance can enable impersonation and monetary loss, while a false rejection
can exclude a legitimate customer or repeatedly expose their biometric data.
The system must therefore balance security, usability, latency, resource use,
privacy, and equitable performance instead of optimising only dataset accuracy.

Deepfake attacks are difficult for several connected reasons:

\begin{itemize}
  \item manipulation mechanisms are heterogeneous and continue to evolve;
  \item compression, resizing, recapture, and post-processing can erase the
        artefacts on which a detector relies;
  \item a model can learn dataset, identity, background, or codec shortcuts and
        then fail under a new generator or population;
  \item a manipulated face can preserve plausible blinking, pose, and motion, so
        conventional liveness cues alone do not establish authenticity;
  \item a virtual camera or compromised client can bypass assumptions about a
        trusted physical capture path; and
  \item mobile inference restricts model size, memory, energy, and response time.
\end{itemize}

The attack must consequently be described along more than one axis.
\Cref{fig:taxonomy} distinguishes the delivery path from the manipulation family.
The five FaceForensics++-derived experimental fake classes are not five equivalent
meanings of ``deepfake'': Deepfakes, FaceSwap, and FaceShifter are identity-swap
approaches, whereas Face2Face and NeuralTextures are reenactment/neural-rendering
approaches \cite{rossler2019faceforensics,li2020faceshifter}. ``Original'' denotes
genuine data and is not a manipulation class.

\begin{figure}[H]
  \centering
  \begin{tikzpicture}[
  font=\small,
  box/.style={draw=buetblue,rounded corners,fill=softblue,align=center,minimum height=11mm,text width=37mm},
  leaf/.style={draw=gray!75,rounded corners,fill=white,align=center,minimum height=10mm,text width=37mm},
  arr/.style={-{Latex[length=2.2mm]},thick,draw=buetblue}
]
\node[box] (capture) {Delivery path};
\node[leaf,below left=12mm and 17mm of capture] (physical) {Physical presentation\\print, replay, mask};
\node[leaf,below=12mm of capture] (recapture) {Deepfake recaptured\\through a display};
\node[leaf,below right=12mm and 17mm of capture] (inject) {Digital injection\\virtual camera or API};
\draw[arr] (capture) -- (physical);
\draw[arr] (capture) -- (recapture);
\draw[arr] (capture) -- (inject);

\node[box,below=23mm of recapture] (family) {Manipulation family};
\node[leaf,below left=12mm and 17mm of family] (swap) {Identity swap\\Deepfakes, FaceSwap,\\FaceShifter};
\node[leaf,below=12mm of family] (reenact) {Reenactment\\Face2Face,\\NeuralTextures};
\node[leaf,below right=12mm and 17mm of family] (outside) {Outside experiment\\full synthesis, lip-sync,\\audio or document fake};
\draw[arr] (family) -- (swap);
\draw[arr] (family) -- (reenact);
\draw[arr] (family) -- (outside);
\end{tikzpicture}


  \caption{Threat taxonomy used in this report. Delivery path and manipulation
  family are separate: the same generated video may be shown to a camera or
  injected directly. Items outside the experiment remain relevant threats.}
  \label{fig:taxonomy}
\end{figure}

\section{Problem Statement}

The engineering problem is to design and evaluate a practical face-liveness
screening pipeline for remote financial verification that detects both physical
presentation attacks and digitally manipulated faces, uses temporal evidence
where useful, degrades safely under uncertainty, and remains feasible on a
mobile--server architecture. The work asks four linked questions:

\begin{enumerate}
  \item Can a lightweight temporal model provide a useful first decision on a
        mobile device?
  \item How well do stronger temporal and domain-generalised PAD models transfer
        from physical-spoof datasets to deepfake datasets?
  \item Does manipulation-specific expertise improve frame-level deepfake
        detection relative to one universal detector?
  \item What security and ethical controls are required before experimental
        accuracy can become a trustworthy financial decision?
\end{enumerate}

For a captured clip $V=(x_1,\ldots,x_K)$, the first-layer model produces a fake
probability $p_1=p(\mathrm{fake}\mid V)$. A target cascade uses two policy
thresholds $\tau_{\mathrm{pass}}<\tau_{\mathrm{attack}}$: clear low-risk cases
may pass, clear high-risk cases may be rejected or challenged, and the interval
between them is escalated. These thresholds are a design requirement; the current
Flutter prototype applies a single threshold of 0.5.

For the later specialist experiment, let $x$ be a preprocessed face frame,
$z=\phi(x)$ an embedding, and $\mathcal{E}$ the five manipulation-specific binary
experts. The intended router is

\begin{align}
q_e(z) &= P(E^*=e\mid z), \qquad e\in\mathcal{E},\\
\hat e &= \arg\max_{e\in\mathcal{E}} q_e(z),\\
\widehat p(\mathrm{fake}\mid x) &= s_{\hat e}(x),
\end{align}

where $s_e$ is the fake probability from expert $e$. The router does not predict
``original''; every frame is assigned to a binary specialist. A truthful claim of
$P(\text{best expert}\mid z)$ additionally requires leakage-safe targets---for
example, choosing the expert with the lowest out-of-fold loss for each training
sample under a declared tie rule. If the training target is merely the known
manipulation label, the correct interpretation is
$P(\text{manipulation method}\mid z)$, not the best expert. The current notebooks
use a six-way original-plus-method classifier and an original-route fallback.
They are therefore reported as a legacy six-way routing result and not as a
completed implementation of the intended equations.

\section{Complexity of the Problem}

The project qualifies as a complex engineering problem because its requirements
conflict and its operating conditions are only partially observable. Important
dimensions are summarised in \Cref{tab:complexity}.

\begin{table}[H]
\centering
\caption{Sources of engineering complexity.}
\label{tab:complexity}
\begin{tabularx}{\textwidth}{@{}p{32mm}Y Y@{}}
\toprule
Dimension & Conflict or uncertainty & Engineering implication \\
\midrule
Attack diversity & Prints, replay, reenactment, face swap, recapture, and
injection expose different signals. & No single benchmark or cue demonstrates
complete protection. \\
Domain shift & Dataset, camera, subject, codec, and generator differ from the
deployment domain. & Group-disjoint and cross-domain evaluation are mandatory. \\
Edge constraints & Larger temporal models may improve evidence but increase
latency, memory, energy, and network use. & Use a lightweight local stage and
selective escalation. \\
Decision cost & False accepts create fraud risk; false rejects create exclusion
and support cost. & Calibrate thresholds against operational costs and provide
retry/appeal paths. \\
Biometric sensitivity & Face data is difficult to revoke and may reveal identity
or demographic attributes. & Minimise collection, secure transit and storage,
and define retention and access controls. \\
Adversarial evolution & Attackers can adapt after learning detector behaviour. &
Monitor drift, limit disclosure of exploitable details, and retest continuously. \\
\bottomrule
\end{tabularx}
\end{table}

\section{Project Objectives and Scope}

The project objectives are to:

\begin{enumerate}
  \item establish a clear physical-presentation, manipulation, and injection
        threat model for remote face verification;
  \item build a lightweight temporal first-stage classifier and integrate its
        exported model into a Flutter capture flow;
  \item investigate richer temporal evidence through $G^2V^2$former;
  \item compare temporal modelling with the GD-FAS domain-generalisation method;
  \item measure within-domain and cross-domain behaviour on public PAD and
        deepfake datasets;
  \item compare a universal Xception detector with manipulation specialists and
        learned routing;
  \item specify a layered escalation architecture and the controls needed for
        trustworthy deployment; and
  \item analyse privacy, fairness, accessibility, accountability, intellectual
        property, cost, and environmental impact.
\end{enumerate}

The implemented scope comprises public-data experiments, preprocessing and
training notebooks, a Flutter/mobile temporal prototype, a separately demonstrated
$G^2V^2$former API, GD-FAS evaluation, and frame-level Xception routing experiments.
It does not comprise a production Pathao Pay deployment, a connected multi-model
backend, capture-device attestation, a reviewer console, calibrated production
thresholds, or a completed field trial. Audio deepfakes, document forgery, fully
synthetic identities, and end-to-end identity matching are threat context but not
separately evaluated model classes.

